\documentclass[krantz1]{krantz}
\usepackage{fixltx2e,fix-cm}
\usepackage{amssymb, bbm, bm, mathtools}
\usepackage{amsmath}
\usepackage[pdftex]{graphicx}
\usepackage{color, rotating}
\usepackage{subcaption}
\usepackage{hyperref}
\usepackage{cleveref}

% Theorem-like environments
\newtheorem{exercise}{Exercise}[chapter]
\newtheorem{example}{Example}
\newtheorem{definition}{Definition}
\newtheorem{assumption}{Assumption}[chapter]
\newtheorem{thesis}{Theorem}[chapter]

% Custom macros from the book
\def\iidsim{\stackrel{\textup{IID}}{\sim}}
\def\indsim{\stackrel{ind}{\sim}}
\def\convergeas{\stackrel{a.s.}{\longrightarrow}}
\def\converged{\stackrel{\textup{d}}{\longrightarrow}}
\def\pr{\textup{pr}}
\def\cov{\textup{cov}}
\def\ind{\perp}
\def\RD{\textsf{rd}}
\def\RR{\textsf{rr}}
\def\OR{\textsf{or}}
\def\true{\textup{true}}
\def\obs{\textup{obs}}
\def\se{\textup{se}}
\def\logit{\textup{logit}}
\def\expit{\textup{expit}}
\def\sumn{\sum_{i=1}^n}
\def\diff{\textsf{d}}
\def\Paragraph{\paragraph}

\begin{document}

\chapter{Problem Set 4}

\section{Book Exercises from Peng Ding's \textit{A First Course in Causal Inference}}

\subsection{Chapter 13: The Average Causal Effect on the Treated Units and Other Estimands}


\begin{definition}[CATE, ATT, and ATE]\label{13:def:CATE-ATT}
The conditional average treatment effect (CATE) is $\tau(X) = E\{ Y(1)-Y(0) \mid X \}$.
The average treatment effect on the treated (ATT) is $\tau_\mathrm{T} = E\{ Y(1)-Y(0) \mid Z=1 \}$.
The average treatment effect on the control (ATC) is $\tau_\mathrm{C} = E\{ Y(1)-Y(0) \mid Z=0 \}$.
The average treatment effect (ATE) is $\tau = E\{ Y(1)-Y(0) \}$.
\end{definition}
 \Paragraph{13.1~Comparing $\tau_\mathrm{T}, \tau_\mathrm{C}$, and $\tau$}\label{hw::difference-tau-T-tau}

Assume $Z\ind \{Y(1), Y(0)\} \mid X$. Recall $e(X) = \pr(Z=1\mid X)$ is the propensity score, $e = \pr(Z=1)$ is the marginal probability of the treatment, and $\tau(X) = E\{Y(1)-Y(0) \mid X\}$ is the CATE.
Show that
$$
\tau_\mathrm{T} - \tau = \frac{\cov\{e(X), \tau(X)\}}{e},\quad
\tau_\mathrm{C} - \tau = - \frac{\cov\{e(X), \tau(X)\}}{1-e}.
$$

Remark: The results above also imply that $\tau_\mathrm{T} - \tau_\mathrm{C} = \cov\{e(X), \tau(X)\} / \{e(1-e)\}$. So the differences among $\tau_\mathrm{T}, \tau_\mathrm{C}, \tau$ depend on the covariance between the propensity score and CATE. When $e(X)$ and $\tau(X)$ are uncorrelated, their differences are 0.

 \Paragraph{13.6~Estimating individual effect and conditional average causal effect}\label{hw::individual-effect-prediction}

Assume that $\{Z_i, X_i, Y_i(1), Y_i(0)\}_{i=1}^n \iidsim \{Z, X, Y(1), Y(0)\}$. The individual effect is $\tau_i = Y_i(1) - Y_i(0)$ and the conditional average causal effect is $\tau(X_i) = E\{Y_i(1) - Y_i(0) \mid X_i\}$.

\begin{enumerate}
\item
Under randomization with $Z_i\ind \{Y_i(1), Y_i(0)\}$ and $e = \pr(Z_i = 1)$, show that
$$
\delta_i = \frac{Z_i Y_i}{e} - \frac{(1-Z_i)Y_i}{1-e}
$$
is an unbiased predictor of the individual effect in the sense that
$$
E(\delta_i - \tau_i) = 0\quad (i=1,\ldots, n).
$$
Further show that $E(\delta_i) = \tau$ for all $i=1,\ldots, n$.

\item
Under ignorability with $Z_i\ind \{Y_i(1), Y_i(0)\}\mid X_i$ and $e(X_i) = \pr(Z_i=1\mid X_i)$, show that
$$
\delta_i = \frac{Z_i Y_i}{e(X_i)} - \frac{(1-Z_i)Y_i}{1-e(X_i)}
$$
is an unbiased predictor of the individual effect and the conditional average causal effect in the sense that
$$
E(\delta_i - \tau_i) = 0, \quad
E\{\delta_i - \tau(X_i)\} = 0, \quad (i=1,\ldots, n).
$$
Further show that $E(\delta_i) = \tau$ for all $i=1,\ldots, n$.
\end{enumerate}


\begin{definition}[General estimand $\tau^h$]\label{13:def:tau-h}
For a given function $h(X)$, the general estimand is
$$
\tau^h = \frac{E\{h(X)\tau(X)\}}{E\{h(X)\}} = \frac{E[h(X)\{\mu_1(X) - \mu_0(X)\}]}{E\{h(X)\}},
$$
where $\mu_1(X) = E\{Y(1)\mid X\}$ and $\mu_0(X) = E\{Y(0)\mid X\}$.
\end{definition}

\begin{definition}[Doubly robust estimator]\label{13:def:DR-estimator}
The doubly robust estimator uses outcome regression models $\mu_1(X,\beta_1)$ and $\mu_0(X,\beta_0)$, and propensity score model $e(X,\alpha)$:
\begin{align*}
\tilde{\mu}_1^{h,\mathrm{dr}} &= E\left[\frac{Zh(X)\{Y-\mu_1(X,\beta_1)\}}{e(X,\alpha)} + h(X)\mu_1(X,\beta_1)\right], \\
\tilde{\mu}_0^{h,\mathrm{dr}} &= E\left[\frac{(1-Z)h(X)\{Y-\mu_0(X,\beta_0)\}}{1-e(X,\alpha)} + h(X)\mu_0(X,\beta_0)\right].
\end{align*}
\end{definition}
 \Paragraph{13.10~Doubly robust estimation for general estimand}\label{hw::general-estimand-dr}

For a given $h(X)$, we have the following formulas for constructing the doubly robust estimator for $\tau^h$:
\begin{align*}
\tilde{\mu}_1^{h,\mathrm{dr}} &= E\left[\frac{Zh(X)\{ Y-\mu_1(X,\beta_1) \}}{e(X,\alpha)} + h(X)\mu_1(X,\beta_1)\right], \\
\tilde{\mu}_0^{h,\mathrm{dr}} &= E\left[\frac{(1-Z)h(X)\{ Y-\mu_0(X,\beta_0) \}}{1-e(X,\alpha)} + h(X)\mu_0(X,\beta_0)\right].
\end{align*}
Show that under ignorability and overlap:
\begin{enumerate}
\item
if either $e(X,\alpha) = e(X)$ or $\mu_1(X,\beta_1) = \mu_1(X)$, then $\tilde{\mu}_1^{h,\mathrm{dr}} = E\{h(X)Y(1)\}$;
\item
if either $e(X,\alpha) = e(X)$ or $\mu_0(X,\beta_0) = \mu_0(X)$, then $\tilde{\mu}_0^{h,\mathrm{dr}} = E\{h(X)Y(0)\}$;
\item
if either $e(X,\alpha) = e(X)$ or $\{\mu_1(X,\beta_1) = \mu_1(X), \mu_0(X,\beta_0) = \mu_0(X)\}$, then
$$
\frac{\tilde{\mu}_1^{h,\mathrm{dr}} - \tilde{\mu}_0^{h,\mathrm{dr}}}{E\{h(X)\}} = \tau^h.
$$
\end{enumerate}

\subsection{Chapter 14: Using the Propensity Score in Regressions for Causal Effects}

 \Paragraph{14.3~Weighted logistic regression with a binary outcome}\label{hw::logistic-regression}

With a binary outcome, we can replace linear outcome models with logistic outcome models. Show that with weights in the logistic regressions, the doubly robust estimators equal the outcome regression estimator. The result holds for both $\tau$ and $\tau_\mathrm{T}$.

\subsection{Chapter 17: E-Value: Evidence for Causation in Observational Studies with Unmeasured Confounding}

\begin{definition}[Risk ratios and risk difference]\label{17:def:RR-RD}
For binary $Y$, the true conditional causal risk ratio is
$$\RR_{ZY\mid x}^\true = \frac{\pr\{Y(1)=1\mid X=x\}}{\pr\{Y(0)=1\mid X=x\}}.$$
The observed conditional risk ratio is
$$\RR_{ZY\mid x}^\obs = \frac{\pr(Y=1\mid Z=1,X=x)}{\pr(Y=1\mid Z=0,X=x)}.$$
The treatment-confounder risk ratio is
$$\RR_{ZU\mid x} = \frac{\pr(U=1\mid Z=1,X=x)}{\pr(U=1\mid Z=0,X=x)}.$$
The confounder-outcome risk ratio is
$$\RR_{UY\mid x} = \frac{\pr(Y=1\mid U=1,X=x)}{\pr(Y=1\mid U=0,X=x)}.$$
The risk difference is $\RD_{ZY} = \pr(Y=1\mid Z=1) - \pr(Y=1\mid Z=0)$.
\end{definition}

\begin{thesis}\label{17.1}\emph{(Theorem 17.1, Ding, 2025)}
Under $Z\ind Y\mid (X, U)$, assume
\begin{equation}\label{eq::evalue-17.1}
\RR_{ZY\mid x}^\obs > 1, \quad \RR_{ZU\mid x} > 1,\quad \RR_{UY\mid x} > 1.
\end{equation}
We have
$$
\RR_{ZY\mid x}^\obs \leq \frac{ \RR_{ZU\mid x}  \RR_{UY\mid x} }{  \RR_{ZU\mid x} + \RR_{UY\mid x} - 1} .
$$
\end{thesis}

 \Paragraph{17.1~A technical lemma for the proof of \cref{17.1}}\label{hw::technical-lemma}

Show that $f(x) = \frac{ax+1}{bx+1}$ is increasing in $x$ if $a > b$ and decreasing in $x$ if $a < b$.

 \Paragraph{17.2~Technical assumption for Theorem 17.1}\label{hw::technical-assumption}

Revisit the proof of \cref{17.1}. Assume $Z\ind Y \mid (X, U)$. Show that if $\RR_{ZU|x} > 1$ but $\RR_{UY|x} < 1$, then $\RR_{ZY|x}^{\obs} < 1$.

Remark: This result is intuitive. Condition on $X$. It says that if the $Z$-$U$ relationship is positive and the $U$-$Y$ relationship is negative, then the $Z$-$Y$ relationship is negative given the conditional independence of $Z$ and $Y$ given $U$.


\begin{definition}[Cornfield-type risk differences]\label{17:def:cornfield-rd}
For binary $Z, Y, U$, the observed risk difference is
$$\RD_{ZY}^{\obs} = \pr(Y=1\mid Z=1) - \pr(Y=1\mid Z=0).$$
The treatment-confounder risk difference is
$$\RD_{ZU} = \pr(U=1\mid Z=1) - \pr(U=1\mid Z=0).$$
The confounder-outcome risk difference is
$$\RD_{UY} = \pr(Y=1\mid U=1) - \pr(Y=1\mid U=0).$$
Under latent ignorability $Z\ind Y\mid U$, we have the product relation $\RD_{ZY}^{\obs} = \RD_{ZU} \times \RD_{UY}$.
\end{definition}

 \Paragraph{17.6~Cornfield-type inequalities for the risk difference}\label{hw::cornfield-rd}

Consider binary $Z, Y, U$, and condition on $X$ implicitly. Assume latent ignorability given $U$. Show that under $Z\ind Y \mid U$, we have
$$
\RD_{ZY}^{\obs} = \RD_{ZU} \times \RD_{UY}
$$
where $\RD_{ZY}^{\obs}$ is the observed risk difference of $Z$ on $Y$, and $\RD_{ZU}$ and $\RD_{UY}$ are the treatment-confounder and confounder-outcome risk differences, respectively.

Remark: Without loss of generality, assume that $\RD_{ZY}^{\obs}, \RD_{ZU}, \RD_{UY}$ are all positive. Then the above implies:
$$
\min(\RD_{ZU}, \RD_{UY}) \geq \RD_{ZY}^{\obs}
$$
and
$$
\max(\RD_{ZU}, \RD_{UY}) \geq \sqrt{\RD_{ZY}^{\obs}}.
$$

These are the Cornfield inequalities for the risk difference with a binary confounder.

\subsection{Chapter 20: Overlap in Observational Studies: Difficulties and Opportunities}

\begin{definition}[Overlap and propensity score]\label{20:def:overlap}
The propensity score is $e(X) = \pr(Z=1\mid X)$.
The overlap condition requires $0 < e(X) < 1$.
The marginal probability of treatment is $e = \pr(Z=1)$.
Let $\lambda_1$ and $\lambda_0$ be the maximum eigenvalues of $\cov(X\mid Z=1)$ and $\cov(X\mid Z=0)$, respectively.
The strict overlap condition requires $\eta \leq e(X) \leq 1-\eta$ for some $\eta \in (0,1/2)$.
\end{definition}


\begin{assumption}[strict overlap]\label{20.1}
$\eta \leq e(X) \leq 1-\eta$ for some $\eta \in (0,1/2).$
\end{assumption}

\begin{thesis}\label{20.2}
Assumption \ref{20.1} implies that $\eta \leq e \leq 1-\eta$ and
\begin{align}
& p^{-1} \sum_{k=1}^p \big | E(X_k\mid Z=1) - E(X_k\mid Z=0) \big | \leq p^{-1/2} C^{1/2} \{ e\lambda_1^{1/2} + (1-e) \lambda_0^{1/2} \},
\end{align}
where
$$
C = \frac{ (e-\eta)(1-e-\eta) }{ e^2(1-e)^2 \eta (1-\eta) }
$$
is a positive constant depending only on $(e, \eta)$, and $\lambda_1$ and $\lambda_0$ are the maximum eigenvalues of the covariance matrices $\cov(X\mid Z=1)$ and $\cov(X\mid Z=0)$, respectively.
\end{thesis}


\beginemark{Relevant computation formulas:}
\begin{enumerate}
\item \textbf{E-value bounding formula:} The upper bound for the observed risk ratio is
$$\RR_{ZY\mid x}^\obs \leq \frac{\RR_{ZU\mid x} \cdot \RR_{UY\mid x}}{\RR_{ZU\mid x} + \RR_{UY\mid x} - 1}.$$
When $\RR_{ZU\mid x} = \RR_{UY\mid x} = r$, the bound simplifies to $r$.
\item \textbf{Constant $C$ in overlap inequality:}
$$C = \frac{(e-\eta)(1-e-\eta)}{e^2(1-e)^2\eta(1-\eta)}.$$
For practical computation, note that $C \approx 1/(e(1-e))$ when $\eta$ is small.
\item \textbf{IPW estimator for general estimand:} $\tau^h = E\{w(X) \cdot ZY / e(X)\} / E\{h(X)\}$ where $w(X) = h(X) + \{1-e(X)\}/e(X) \cdot something$.
\item \textbf{Risk difference product formula:} $\RD_{ZY}^\obs = \RD_{ZU} \times \RD_{UY}$ for binary confounders.
\end{enumerate}
\endremark}

 \Paragraph{20.1~Implications of overlap}\label{hw::overlap-implications}

In Part III of this book, causal inference with observational studies relies on two critical assumptions: ignorability $Z\ind \{Y(1), Y(0)\} \mid X$ and overlap $0 < e(X) < 1$.

D'Amour et al. (2021) pointed out the tension between these two assumptions: typically, more covariates make the ignorability assumption more plausible, but more covariates make the overlap assumption less plausible because the treatment becomes more predictable given more covariates.

Many statistical analyses require a strict version of overlap (Assumption \ref{20.1}).

Show that Assumption \ref{20.1} implies $\eta \leq e \leq 1 - \eta$ and the inequality above.

\section{Theoretical Exercises: Identification Formulas for Sensitivity Analysis}

Denote $\mu_z(X) = E(Y \mid Z=z, X)$ and $e(X) = \pr(Z=1 \mid X)$.
Suppose that
$$
\frac{E\{Y(1) \mid Z=1, X\}}{E\{Y(1) \mid Z=0, X\}} = \epsilon_1(X), \quad
\frac{E\{Y(0) \mid Z=1, X\}}{E\{Y(0) \mid Z=0, X\}} = \epsilon_0(X)
$$
with known $\epsilon_1(X)$ and $\epsilon_0(X)$.

 \Paragraph{(a)~Identification formulas}\label{hw::sensitivity-id}

Show that
$$
E\{Y(1) \mid Z=0\} = E\{\mu_1(X)/\epsilon_1(X) \mid Z=0\}, \quad
E\{Y(0) \mid Z=1\} = E\{\mu_0(X)\epsilon_0(X) \mid Z=1\}
$$
and hence
$$
\mathrm{ACE} = E\{ZY + (1-Z)\mu_1(X)/\epsilon_1(X)\} - E\{Z\mu_0(X)\epsilon_0(X) + (1-Z)Y\}.
$$

 \Paragraph{(b)~IPW representation}\label{hw::sensitivity-ipw}

Show that
$$
E\{Y(1)\} = E\left\{w_1(X)\frac{ZY}{e(X)}\right\}, \quad
E\{Y(0)\} = E\left\{w_0(X)\frac{(1-Z)Y}{1-e(X)}\right\},
$$
where $w_1(X) = e(X) + \{1-e(X)\}/\epsilon_1(X)$ and $w_0(X) = e(X)\epsilon_0(X) + \{1-e(X)\}$.

 \Paragraph{(c)~Doubly robust estimation (Optional)}\label{hw::sensitivity-dr}

Derive the doubly robust estimation formula.

\end{document}
